\documentclass[12pt]{article}

% Basic packages
\usepackage{geometry}
\usepackage{amsfonts}
\usepackage{amsmath}
\usepackage{amssymb}
\usepackage{graphicx}
\usepackage[small,bf,center]{caption}
\usepackage{subcaption}
\usepackage{setspace}
\usepackage{float}
\usepackage[autostyle]{csquotes}

% Packages for automated tables
\usepackage{booktabs}
\usepackage{tabularx}
\usepackage{standalone}
\usepackage{pdflscape}

\onehalfspacing
\textwidth=6.0in
\textheight=8.5in
\begin{document}

\title{My Paper}
\author{Julian Reif\thanks{University of Illinois and NBER.}}
\maketitle

\begin{abstract}
\noindent 
This paper provides an example of a document with tables and figures that were automated using Stata.
\end{abstract}

\clearpage

\section{Summary}

This example paper includes a table and a figure produced by \texttt{scripts/1\_example.do}, which writes them directly into \texttt{paper/figures/} and \texttt{paper/tables/} so this document compiles after a fresh run of \texttt{run.do}.

I use Stata's \textbf{auto.dta} dataset.%
%
\footnote{Type \textbf{sysuse auto, clear} at the Stata prompt to load this dataset.}
%
The price distribution, illustrated in Figure \ref{fig:price_histogram}, is skewed right.

I estimate the association between automobile prices and fuel efficiency using the following linear model:
\begin{align}
PRICE_i&=\alpha + \beta X_i + \varepsilon \label{eqn:model}
\end{align}
The outcome variable, $PRICE_i$, is the price of automobile $i$. The parameter of interest is $\beta$, a vector of coefficients. In my first specification (``spec 1''), the vector $X_i$ includes miles per gallon. The second specification (``spec 2'') also includes the car's weight. I estimate this model using ordinary least squares and report standard errors that are robust to heteroskedasticity.

Table \ref{tab:my_regressions} reports my estimates, separately for domestic and foreign cars. Column (1) reports that an increase in fuel efficiency of 1 mile per gallon is associated with a \$329 reduction in the price of domestic automobiles. Column (2) shows that this association becomes positive and insignificant when I also include weight as a regressor. Columns (3) and (4) show that these associations are similar for foreign automobiles.

\clearpage
\section{Figures and Tables}


%%%
% Figure: Price histogram
%%%
\begin{figure}[ht]
\caption{Automobile prices}\label{fig:price_histogram}
\begin{center}
{\includegraphics[width=1\textwidth]{./figures/price_histogram.png}}
\parbox{\linewidth}{\footnotesize {Notes: Data were obtained from Stata's built-in auto dataset.}}
\end{center}


\end{figure}
\clearpage

%%%
% Tables
%%%

\input{tables/my_regressions.tex}

\end{document}